\chapter{Derivación y teoremas del valor medio}
\label{chap:derivacion-valor-medio}

La derivada mide una variación local. Puede leerse como pendiente de una
tangente, como tasa instantánea o como la mejor aproximación lineal de una
función cerca de un punto. Estas tres lecturas describen el mismo número, pero
la aproximación lineal será la que organice las reglas de cálculo y evite
argumentos formales incompletos.

\section{La derivada como variación local}
\label{sec:derivada-variacion-local}

\subsection{Cociente incremental, pendiente y tasa}
\CanonicalUnit{CM-C03-U001}{}{}

Sea \(f:D\to\R\) y sea \(a\) un punto interior de \(D\). La tasa media de
cambio entre \(a\) y \(a+h\), con \(h\ne0\), es
\[
  \frac{f(a+h)-f(a)}{h}.
\]
Geométricamente es la pendiente de la recta secante que une
\((a,f(a))\) con \((a+h,f(a+h))\).

\begin{definicion}[Derivada]
La función \(f\) es \emph{derivable en \(a\)} si existe y es finito el límite
\[
  f'(a)=\lim_{h\to0}\frac{f(a+h)-f(a)}{h}.
\]
Equivalentemente,
\[
  f'(a)=\lim_{x\to a}\frac{f(x)-f(a)}{x-a}.
\]
El número \(f'(a)\) es la \emph{derivada} de \(f\) en \(a\).
\end{definicion}

La variable de incremento \(h\) tiende a \(0\), no a \(a\). Si la derivada
existe en todos los puntos de un conjunto, la \emph{función derivada} asigna
\(x\mapsto f'(x)\).

Cuando solo hay puntos del dominio a un lado de \(a\), se definen las
derivadas laterales mediante \(h\to0^+\) o \(h\to0^-\). Si ambos lados están
disponibles, la derivada bilateral existe si y solo si las dos derivadas
laterales existen y coinciden.

El paso de secantes a la recta tangente corresponde al límite de sus
pendientes. En un modelo temporal, el mismo cociente es una tasa media y su
límite es una tasa instantánea.

\begin{ejemplo}[Velocidad instantánea]
Si la posición de una partícula es \(s(t)=t^2\), su velocidad media entre
\(t=3\) y \(t=3+h\) es
\[
  \frac{s(3+h)-s(3)}{h}=6+h.
\]
Al hacer \(h\to0\), la velocidad instantánea es \(s'(3)=6\), en unidades de
longitud por unidad de tiempo.
\end{ejemplo}

\subsection{Aproximación lineal y diferencial}
\CanonicalUnit{CM-C03-U002}{}{}

La existencia de \(f'(a)\) equivale a poder escribir
\[
  f(a+h)=f(a)+f'(a)h+r(h),
  \qquad
  \frac{r(h)}{h}\longrightarrow0.
\]
Con la notación del Capítulo~\ref{chap:sucesiones-limites-continuidad},
\[
  f(a+h)=f(a)+f'(a)h+o(h)
  \qquad(h\to0).
\]
El término \(f(a)+f'(a)h\) es la aproximación lineal local; el error es de
orden superior al incremento.

\begin{definicion}[Diferencial en un punto]
La \emph{diferencial de \(f\) en \(a\)} es la aplicación lineal
\[
  df_a:\R\to\R, \qquad df_a(k)=f'(a)k.
\]
\end{definicion}

La diferencial no es otro nombre para el número \(f'(a)\): es la aplicación
lineal determinada por ese número. En una variable ambas informaciones se
identifican fácilmente, pero la distinción aclara la regla de la cadena y
prepara generalizaciones posteriores.

\subsection{Derivabilidad y continuidad}
\CanonicalUnit{CM-C03-U003}{}{}

\begin{teorema}[La derivabilidad implica continuidad]
Si \(f\) es derivable en \(a\), entonces es continua en \(a\).
\end{teorema}

\begin{proof}
La aproximación lineal da
\[
  f(a+h)-f(a)=f'(a)h+r(h),
  \qquad r(h)=o(h).
\]
Cuando \(h\to0\), tanto \(f'(a)h\) como \(r(h)\) tienden a cero. Por tanto,
\(f(a+h)\to f(a)\), que es la continuidad en \(a\).
\end{proof}

El recíproco es falso. La función \(f(x)=|x|\) es continua en \(0\), pero sus
derivadas laterales valen \(-1\) y \(1\), por lo que no es derivable allí. La
función \(x^{1/3}\) también es continua en cero, pero su cociente incremental
es \(h^{-2/3}\), que no tiene límite real finito. Una esquina y una tangente
vertical son dos mecanismos distintos de no derivabilidad.

\section{Derivadas elementales y reglas algebraicas}
\label{sec:reglas-derivacion}

\subsection{Derivaciones representativas}
\CanonicalUnit{CM-C03-U004}{}{}
\ProofRole{proof-018}{}
\ProofRole{proof-004}{}

La derivada de una constante es cero y la de la identidad es uno. Para
\(n\in\N\), el binomio de Newton da
\[
 \frac{(a+h)^n-a^n}{h}
 =\sum_{k=1}^{n}\binom nk a^{n-k}h^{k-1},
\]
de donde
\[
  \frac{d}{dx}x^n=nx^{n-1}.
\]
Las reglas algebraicas extenderán esta fórmula a potencias enteras negativas y
la regla de la función inversa permitirá tratar raíces en sus ramas reales.

Para el número \(e\) definido en el capítulo anterior, la construcción de la
exponencial real proporciona la caracterización equivalente
\[
  \lim_{h\to0}\frac{e^h-1}{h}=1.
\]
Es una propiedad de referencia de esa construcción, no una nueva definición
independiente. Por tanto,
\[
 \frac{e^{x+h}-e^x}{h}
 =e^x\frac{e^h-1}{h}\longrightarrow e^x,
 \qquad
 (e^x)'=e^x.
\]
Si \(b>0\), entonces \(b^x=e^{x\ln b}\), y la regla de la cadena dará
\[
  (b^x)'=b^x\ln b.
\]
La fórmula incluye \(b=1\), cuyo caso es la función constante; no utiliza un
logaritmo de base uno.

Para el seno usamos la identidad de adición:
\[
 \frac{\sin(x+h)-\sin x}{h}
 =\sin x\frac{\cos h-1}{h}
  +\cos x\frac{\sin h}{h}.
\]
Los límites del Capítulo~2 dan
\[
  (\sin x)'=\cos x.
\]
Del mismo modo, \((\cos x)'=-\sin x\). Después de aplicar las reglas de
producto, cociente, cadena e inversa, obtendremos las restantes fórmulas
elementales sin repetir todos los límites desde la definición.

\subsection{Suma, producto y cociente}
\CanonicalUnit{CM-C03-U005}{}{}
\ProofRole{proof-019}{}

\begin{teorema}[Reglas algebraicas de derivación]
Sean \(f\) y \(g\) derivables en \(a\), y sea \(\lambda\in\R\). Entonces
\[
 (f+g)'(a)=f'(a)+g'(a),
 \qquad
 (\lambda f)'(a)=\lambda f'(a),
\]
\[
 (fg)'(a)=f'(a)g(a)+f(a)g'(a).
\]
Si además \(g(a)\ne0\), el cociente está definido en un entorno de \(a\) y
\[
 \left(\frac f g\right)'(a)
 =\frac{f'(a)g(a)-f(a)g'(a)}{g(a)^2}.
\]
\end{teorema}

\begin{proof}
Las reglas de suma y múltiplo escalar se obtienen distribuyendo el cociente
incremental. Para el producto, sumamos y restamos \(f(a+h)g(a)\):
\[
\begin{aligned}
 \frac{f(a+h)g(a+h)-f(a)g(a)}{h}
 &=f(a+h)\frac{g(a+h)-g(a)}h\\
 &\quad+g(a)\frac{f(a+h)-f(a)}h.
\end{aligned}
\]
La derivabilidad implica continuidad, así que \(f(a+h)\to f(a)\), y el límite
es la fórmula del producto.

Para el recíproco, la continuidad de \(g\) y \(g(a)\ne0\) garantizan que
\(g(a+h)\ne0\) cuando \(h\) es pequeño. Entonces
\[
 \frac{1/g(a+h)-1/g(a)}h
 =-\frac{1}{g(a+h)g(a)}
   \frac{g(a+h)-g(a)}h
 \longrightarrow-\frac{g'(a)}{g(a)^2}.
\]
La regla del producto aplicada a \(f(1/g)\) produce la fórmula del cociente.
\end{proof}

En particular,
\[
  (\tan x)'=\frac1{\cos^2x}
\]
en los puntos donde \(\cos x\ne0\). Las condiciones de dominio no desaparecen
al derivar.

\subsection{Derivadas de orden superior}
\CanonicalUnit{CM-C03-U006}{}{}

Si la función derivada \(f'\) es a su vez derivable, definimos la segunda
derivada \(f''=(f')'\). Recursivamente,
\[
  f^{(n)}=(f^{(n-1)})', \qquad n\ge2,
\]
en los puntos donde todas las derivadas anteriores existen. No debe suponerse
que la existencia de \(f'\) implique automáticamente la de \(f''\).

En un modelo de posición \(s(t)\), \(s'(t)\) representa velocidad y
\(s''(t)\) aceleración. Esta interpretación depende de las unidades y del
modelo, no solo de la fórmula simbólica.

\section{Composición, inversa y técnicas de derivación}
\label{sec:composicion-inversa-derivacion}

\subsection{Regla de la cadena}
\CanonicalUnit{CM-C03-U007}{}{}
\ProofRole{proof-020}{}

\begin{teorema}[Regla de la cadena]
Sean \(A,B,C,D\subseteq\R\), \(f:A\to B\) y \(g:C\to D\), con
\(f(A)\subseteq C\). Si \(f\) es derivable en \(a\in A\) y \(g\) es
derivable en \(b=f(a)\), entonces
\(g\circ f\) es derivable en \(a\) y
\[
  (g\circ f)'(a)=g'(f(a))f'(a).
\]
\end{teorema}

\begin{proof}
Por derivabilidad de \(f\), existe un resto \(r_f(h)=o(h)\) tal que
\[
  f(a+h)=b+f'(a)h+r_f(h).
\]
Definimos el incremento interior
\(k(h)=f(a+h)-b\). Entonces \(k(h)\to0\) y el cociente \(k(h)/h\) permanece
acotado cuando \(h\to0\).

Por derivabilidad de \(g\), podemos escribir
\[
  g(b+k)=g(b)+g'(b)k+k\eta(k),
\]
donde \(\eta(k)\to0\) cuando \(k\to0\); definimos \(\eta(0)=0\). Al sustituir
\(k=k(h)\),
\[
\begin{aligned}
 g(f(a+h))-g(f(a))
 &=g'(b)\bigl(f'(a)h+r_f(h)\bigr)\\
 &\quad+k(h)\eta(k(h)).
\end{aligned}
\]
Dividiendo por \(h\), el término con \(r_f(h)\) tiende a cero. El último
también, porque \(k(h)/h\) está acotado y \(\eta(k(h))\to0\). Queda
\(g'(b)f'(a)\). El argumento incluye sin cancelación los valores de \(h\) para
los que \(k(h)=0\).
\end{proof}

En lenguaje de diferenciales,
\[
  d(g\circ f)_a=dg_{f(a)}\circ df_a.
\]
La regla expresa que las aproximaciones lineales se componen en el mismo orden
que las funciones.

\subsection{Derivada de la función inversa}
\CanonicalUnit{CM-C03-U008}{}{}
\ProofRole{proof-021}{}

\begin{teorema}[Derivada de la función inversa]
Sea \(f:I\to J\) una función biyectiva entre intervalos, con
\(x_0\in\operatorname{int}I\) y \(y_0=f(x_0)\in\operatorname{int}J\).
Supongamos que \(f^{-1}\) es continua en \(y_0\), que \(f\) es derivable en
\(x_0\) y que \(f'(x_0)\ne0\). Entonces \(f^{-1}\) es
derivable en \(y_0\) y
\[
  (f^{-1})'(y_0)=\frac1{f'(x_0)}.
\]
\end{teorema}

\begin{proof}[Estrategia no circular]
Para \(y\ne y_0\), ponemos \(x=f^{-1}(y)\). La inyectividad garantiza
\(x\ne x_0\), y
\[
 \frac{f^{-1}(y)-f^{-1}(y_0)}{y-y_0}
 =\frac{x-x_0}{f(x)-f(x_0)}.
\]
Cuando \(y\to y_0\), la continuidad de \(f^{-1}\) implica \(x\to x_0\).
El último cociente tiende a \(1/f'(x_0)\), que existe porque
\(f'(x_0)\ne0\). Así se establece primero la derivabilidad de la inversa; la
regla de la cadena puede comprobar después la fórmula, pero no se usa para
suponer la conclusión.
\end{proof}

La exponencial natural es biyectiva de \(\R\) en \((0,+\infty)\). Aplicando el
teorema a \(f(x)=e^x\), con \(x_0=\ln y_0\), obtenemos
\[
  (\ln y)'=\frac1y, \qquad y>0.
\]

Para \(n\) impar, \(x\mapsto x^n\) es biyectiva de \(\R\) en \(\R\). Para
\(n\) par debe restringirse, por ejemplo, a
\([0,+\infty)\to[0,+\infty)\) antes de invertirla como raíz \(n\)-ésima
real. En puntos con raíz no nula, la fórmula inversa da
\[
  \frac{d}{dx}x^{1/n}
  =\frac{1}{n\bigl(x^{1/n}\bigr)^{n-1}}
\]
en la rama y dominio declarados.

Del mismo modo, usando las ramas principales del Capítulo~1,
\[
 \frac{d}{dx}\operatorname{arcsen}x=\frac1{\sqrt{1-x^2}}
 \quad(|x|<1),
\]
\[
 \frac{d}{dx}\arccos x=-\frac1{\sqrt{1-x^2}}
 \quad(|x|<1),
 \qquad
 \frac{d}{dx}\arctan x=\frac1{1+x^2}
 \quad(x\in\R).
\]
Aquí \(\operatorname{arcsen}\) es el arcoseno; \(\arcsin\) es el alias
internacional usado en el capítulo de funciones.

\subsection{Derivación implícita}
\CanonicalUnit{CM-C03-U009}{}{}

Una relación entre \(x\) e \(y\) puede describir localmente \(y\) como función
de \(x\) sin que esa función esté despejada. La derivación implícita parte de
una hipótesis explícita: cerca del punto estudiado existe una función derivable
\(y=y(x)\) que satisface la relación.

En la circunferencia
\[
  x^2+y^2=1,
\]
suponemos que una rama se escribe localmente como \(y=y(x)\). Al derivar la
identidad \(x^2+y(x)^2=1\) y aplicar la regla de la cadena,
\[
  2x+2y(x)y'(x)=0.
\]
Si \(y(x)\ne0\), entonces
\[
  y'(x)=-\frac{x}{y(x)}.
\]
En los puntos con \(y=0\) no puede dividirse por \(2y\); precisamente allí la
circunferencia no se representa localmente como una función derivable de
\(x\) con pendiente finita.

De forma general, si una relación \(F(x,y)=0\) admite localmente una función
derivable \(y(x)\), la regla de la cadena conduce formalmente a
\[
  F_x(x,y)+F_y(x,y)y'(x)=0.
\]
Cuando las derivadas indicadas existen y \(F_y(x,y)\ne0\), puede aislarse
\(y'=-F_x/F_y\). Esta regla calcula la derivada bajo la hipótesis de función
local; no es una demostración del teorema general de la función implícita ni
garantiza por sí sola que dicha función exista.

\subsection{Derivación logarítmica}
\CanonicalUnit{CM-C03-U010}{}{}
\ProofRole{proof-022}{}

Si \(y=f(x)>0\) en un intervalo y \(f\) es derivable, se puede tomar el
logaritmo natural y derivar
\[
  \ln y=\ln f(x),
  \qquad
  \frac{y'}y=\frac{f'(x)}{f(x)}.
\]
La positividad no es opcional: garantiza que el logaritmo real esté definido y
que la división por \(y\) sea válida.

\begin{ejemplo}
Para \(x>0\), definimos \(y=x^x>0\). Entonces
\[
  \ln y=x\ln x.
\]
Al derivar,
\[
  \frac{y'}y=\ln x+1,
  \qquad
  y'=x^x(\ln x+1).
\]
El dominio \(x>0\) forma parte del resultado real.
\end{ejemplo}

\section{De los extremos locales al valor medio}
\label{sec:rolle-valor-medio}

\subsection{Condición necesaria de extremo interior}
\CanonicalUnit{CM-C03-U011}{}{}

\begin{teorema}[Teorema de Fermat]
Si \(f\) tiene un máximo o mínimo local en un punto interior \(c\) de su
dominio y es derivable en \(c\), entonces \(f'(c)=0\).
\end{teorema}

\begin{proof}
Supongamos que \(c\) es un máximo local. Para \(h>0\) pequeño,
\[
  \frac{f(c+h)-f(c)}h\le0,
\]
mientras que para \(h<0\) pequeño el mismo cociente es mayor o igual que cero,
porque el numerador sigue siendo no positivo y el denominador cambia de signo.
Si el límite bilateral existe, debe ser simultáneamente menor o igual y mayor
o igual que cero; por tanto, \(f'(c)=0\). El caso de un mínimo es análogo.
\end{proof}

La conclusión es necesaria, no suficiente: \(f'(c)=0\) no garantiza por sí
solo un extremo.

\subsection{Teorema de Rolle}
\CanonicalUnit{CM-C03-U012}{}{}
\ProofRole{proof-023}{}
\ProofRole{proof-015}{}

\begin{teorema}[Teorema de Rolle]
Sea \(f:[a,b]\to\R\), con \(a<b\). Si \(f\) es continua en \([a,b]\),
derivable en \((a,b)\) y \(f(a)=f(b)\), entonces existe \(c\in(a,b)\) tal
que \(f'(c)=0\).
\end{teorema}

\begin{proof}
Por el teorema del valor extremo, \(f\) alcanza máximo y mínimo absolutos en
\([a,b]\). Si ambos valores coinciden, \(f\) es constante y cualquier
\(c\in(a,b)\) satisface \(f'(c)=0\). Si son distintos, sea
\(q=f(a)=f(b)\). Entonces el máximo es mayor que \(q\) o el mínimo es menor
que \(q\); el extremo correspondiente no puede alcanzarse solamente en
\(a\) o en \(b\), y por tanto se alcanza en algún punto interior. En ese
punto, el teorema de Fermat da \(f'(c)=0\).
\end{proof}

Cada hipótesis tiene un papel: la continuidad cerrada garantiza extremos, la
derivabilidad interior permite aplicar Fermat y la igualdad de valores en los
extremos fuerza un extremo interior salvo en el caso constante.

\subsection{Teorema del valor medio}
\CanonicalUnit{CM-C03-U013}{}{}
\ProofRole{proof-024}{}

\begin{teorema}[Teorema del valor medio]
Sea \(f:[a,b]\to\R\), con \(a<b\). Si \(f\) es continua en \([a,b]\) y
derivable en \((a,b)\), entonces existe \(c\in(a,b)\) tal que
\[
  f'(c)=\frac{f(b)-f(a)}{b-a}.
\]
\end{teorema}

\begin{proof}
Sea
\[
  m=\frac{f(b)-f(a)}{b-a}
\]
y definamos \(\varphi(x)=f(x)-f(a)-m(x-a)\). La función \(\varphi\) es
continua en \([a,b]\), derivable en \((a,b)\), y
\(\varphi(a)=\varphi(b)=0\). Por Rolle existe \(c\in(a,b)\) tal que
\[
  0=\varphi'(c)=f'(c)-m.
\]
De aquí resulta la fórmula.
\end{proof}

El teorema relaciona una tasa media global con una tasa instantánea interior.
No afirma que el punto \(c\) sea único.

\section{Regla de l'Hôpital}
\label{sec:lhopital}

\subsection{Enunciado y control de hipótesis}
\CanonicalUnit{CM-C03-U014}{}{}
\ProofRole{proof-025}{}

La regla de l'Hôpital se estudia después de Rolle y del teorema del valor
medio porque su justificación depende de resultados de valor medio. Es un
teorema para ciertas formas indeterminadas, no una instrucción automática para
derivar numerador y denominador.

\begin{teorema}[Regla de l'Hôpital, forma lateral]
Sean \(f\) y \(g\) derivables en \((a,a+\eta)\), con \(\eta>0\), y supongamos
que
\[
  g(x)\ne0, \qquad g'(x)\ne0
  \quad\text{para }a<x<a+\eta.
\]
Supongamos además que se cumple una de las dos alternativas:
\[
  f(x)\to0,\quad g(x)\to0
  \qquad(x\to a^+),
\]
o bien
\[
  |f(x)|\to+\infty,\quad |g(x)|\to+\infty
  \qquad(x\to a^+).
\]
Si existe, finito o infinito, el límite
\[
  \lim_{x\to a^+}\frac{f'(x)}{g'(x)}=L,
\]
entonces
\[
  \lim_{x\to a^+}\frac{f(x)}{g(x)}=L.
\]
\end{teorema}

Hay versiones análogas para \(x\to a^-\) y \(x\to\pm\infty\), obtenidas con
los cambios de variable correspondientes. No desarrollaremos la demostración
completa, que usa una forma generalizada del teorema del valor medio. Para
aplicar la regla hay que comprobar, en este orden: dominio lateral, forma
\(0/0\) o \(\infty/\infty\), derivabilidad, no anulación del denominador y de
su derivada, y existencia del límite del cociente de derivadas.

\subsection{Aplicaciones directas}
\CanonicalUnit{CM-C03-U015}{}{}
\ProofRole{proof-025}{}

En el límite
\[
  \lim_{x\to+\infty}\frac{\ln x}{x},
\]
numerador y denominador tienden a \(+\infty\). Para \(x>0\), las hipótesis de
derivación se cumplen y
\[
  \lim_{x\to+\infty}\frac{(\ln x)'}{(x)'}
  =\lim_{x\to+\infty}\frac1x=0.
\]
Por l'Hôpital, el límite original es cero.

Si \(\alpha\in\R\) y \(\beta\ne0\), entonces
\[
 \lim_{x\to0}\frac{\sin(\alpha x)}{\sin(\beta x)}
 =\lim_{x\to0}
   \frac{\alpha\cos(\alpha x)}{\beta\cos(\beta x)}
 =\frac\alpha\beta.
\]
La condición \(\beta\ne0\) asegura que el denominador original no sea la
función nula y que el cociente de derivadas esté definido cerca de cero. Este
límite también se obtiene de \(\sin u/u\to1\), a menudo con menos trabajo.

Antes de repetir l'Hôpital hay que verificar de nuevo la forma. Por ejemplo,
\[
 \lim_{x\to0}\frac{1-\cos x}{x^2}
 =\lim_{x\to0}\frac{\sin x}{2x}
 =\lim_{x\to0}\frac{\cos x}{2}
 =\frac12,
\]
y cada paso usa una forma \(0/0\) hasta llegar a un límite directo.

\subsection{Productos y diferencias indeterminadas}
\CanonicalUnit{CM-C03-U016}{}{}
\ProofRole{proof-025}{}

Una forma \(0\cdot\infty\) debe transformarse en un cociente. Para
\(x\to0^+\),
\[
  x\ln x=\frac{\ln x}{1/x}.
\]
El cociente tiene forma \((-\infty)/(+\infty)\), y
\[
  \frac{(\ln x)'}{(1/x)'}
  =\frac{1/x}{-1/x^2}=-x\longrightarrow0.
\]
Por tanto, \(x\ln x\to0\).

Una forma \(\infty-\infty\) suele transformarse mediante denominador común,
factorización o racionalización. Por ejemplo,
\[
 \sqrt{x^2+x}-x
 =\frac{x}{\sqrt{x^2+x}+x}
 =\frac1{\sqrt{1+1/x}+1}
 \longrightarrow\frac12
 \qquad(x\to+\infty).
\]
La racionalización resuelve el límite sin l'Hôpital. Transformar bien la forma
es más importante que forzar una técnica concreta.

\subsection{Formas de potencia}
\CanonicalUnit{CM-C03-U017}{}{}
\ProofRole{proof-025}{}

Para una potencia real variable \(y(x)=f(x)^{g(x)}\) se requiere
\(f(x)>0\) en un entorno lateral del punto. Entonces
\[
  \ln y(x)=g(x)\ln f(x).
\]
La forma de potencia se estudia calculando primero el límite del exponente
logarítmico. Si este límite es \(A\in\R\), la continuidad de la exponencial da
\(y(x)\to e^A\).

Para la forma \(1^\infty\),
\[
 \ln\left(\left(1+\frac1x\right)^x\right)
 =\frac{\ln(1+1/x)}{1/x}
 \longrightarrow1
 \qquad(x\to+\infty),
\]
pues el cociente de derivadas es \(x/(x+1)\). Por tanto,
\[
  \left(1+\frac1x\right)^x\longrightarrow e,
\]
en concordancia con la definición secuencial del capítulo anterior.

La forma \(0^0\) aparece en \(x^x\) cuando \(x\to0^+\). Como
\(x\ln x\to0\), se obtiene \(x^x\to e^0=1\). Para la forma
\(\infty^0\),
\[
  x^{1/x}\longrightarrow1
  \qquad(x\to+\infty),
\]
porque \(\ln x/x\to0\). En todos los casos la positividad de la base y la
forma real deben comprobarse antes de tomar logaritmos.

\begin{advertencia}
La regla de l'Hôpital no es el método por defecto para límites difíciles. Antes
de usarla se identifica la forma, se comprueban las hipótesis y se comparan
métodos alternativos: álgebra, equivalencias, encaje o racionalización pueden
ser más breves y mostrar mejor la estructura del límite. Las comparaciones
mediante Taylor pertenecen a un capítulo posterior.
\end{advertencia}
